%%
%% Speech Analytics and Natural Language Processing Applied to
%% Clinical Communication: A Scoping Review
%%
%% Authors: Rafael Baena Neto, Vicente Idalberto Becerra Sablón,
%%          Fabricio Evangelista dos Santos, Lucas Kenso Terrazzin,
%%          Murilo Minutti
%%
%% Target journal: International Journal of Medical Informatics (IJMI)
%% APC covered: CAPES/Elsevier Read & Publish 2026–2028
%% Class: cas-sc.cls v2.4 (single-column, for submission/review)
%%

\documentclass[a4paper,fleqn]{cas-sc}

\usepackage[authoryear,longnamesfirst]{natbib}
\usepackage[utf8]{inputenc}
\usepackage[T1]{fontenc}
\usepackage{lmodern}
\usepackage{amsmath}
\usepackage{graphicx}
\usepackage{booktabs}
\usepackage{longtable}
\usepackage{array}
\usepackage{hyperref}
\usepackage{xcolor}

\hypersetup{colorlinks=true, linkcolor=blue, citecolor=blue, urlcolor=blue}

%% ================================================================
\begin{document}
\let\WriteBookmarks\relax
\def\floatpagepagefraction{1}
\def\textpagefraction{.001}

%% ----------------------------------------------------------------
%% Running head
%% ----------------------------------------------------------------
\shorttitle{Speech Analytics in Clinical Communication}
\shortauthors{R. Baena Neto et~al.}

%% ----------------------------------------------------------------
%% Title
%% ----------------------------------------------------------------
\title[mode=title]{Speech Analytics and Natural Language Processing
Applied to Clinical Communication: A Scoping Review}

%% ----------------------------------------------------------------
%% Authors and affiliations
%% Affiliation [1]: Graduate Program in Health Data Science, USF
%% Affiliation [2]: Department of Computer Engineering, USF
%% ----------------------------------------------------------------
\author[1]{Rafael Baena Neto}[%
  type=editor,
  orcid=0009-0005-2122-4665,
  role=Researcher
]
\cormark[1]
\fnmark[1]
\ead{rafael.baena@mail.usf.edu.br}
\credit{Conceptualisation, Methodology, Data curation, Formal analysis,
        Writing -- original draft}

\affiliation[1]{%
  organization={Graduate Program in Health Data Science,
                Universidade São Francisco (USF)},
  addressline={Rua Alexandre Rodrigues Barbosa, 45},
  city={Bragança Paulista},
  postcode={12916-900},
  state={SP},
  country={Brazil}
}

\author[1]{Vicente Idalberto Becerra Sablón}[%
  orcid=0000-0003-3127-1906
]
\fnmark[2]
\ead{vicente.sablon@usf.edu.br}
\credit{Supervision, Conceptualisation, Writing -- review and editing}

\author[2]{Fabricio Evangelista dos Santos}[%
  orcid=0000-0000-0000-0000
]
\fnmark[3]
\ead{fabricio.evangelista@mail.usf.edu.br}
\credit{Investigation, Data curation, Methodology, Software,
        Writing -- review and editing}

\author[2]{Lucas Kenso Terrazzin}[%
  orcid=0000-0000-0000-0000
]
\fnmark[4]
\ead{lucas.ida@mail.usf.edu.br}
\credit{Investigation, Data curation, Formal analysis,
        Writing -- review and editing}

\author[2]{Murilo Minutti}[%
  orcid=0000-0000-0000-0000
]
\fnmark[5]
\ead{murilo.minutti@mail.usf.edu.br}
\credit{Investigation, Data curation, Formal analysis, Visualisation,
        Writing -- review and editing}

\affiliation[2]{%
  organization={Department of Computer Engineering,
                Universidade São Francisco (USF)},
  addressline={Rua Alexandre Rodrigues Barbosa, 45},
  city={Bragança Paulista},
  postcode={12916-900},
  state={SP},
  country={Brazil}
}

\cortext[1]{Corresponding author}
\fntext[1]{}
\fntext[2]{}
\fntext[3]{}
\fntext[4]{}
\fntext[5]{}

%% ----------------------------------------------------------------
%% Abstract — max 250 words (IJMI)
%% ----------------------------------------------------------------
\begin{abstract}
\textbf{Background:}
Professional--patient communication is a dense yet predominantly unstructured
and analytically underutilised source of clinical information. Advances in
automatic speech recognition (ASR), natural language processing (NLP), and
vocal biomarker analysis have expanded the capacity to extract structured data
from clinical interactions; nonetheless, few studies operationalise these
interactions as measurable speech-derived variables, particularly
communicational and behavioural markers relevant to patient assessment.

\textbf{Objective:}
To identify, synthesise, and critically evaluate evidence on the application
of speech analytics and NLP to professional--patient clinical communication,
mapping methodological approaches, identified markers, technical challenges,
and research gaps.

\textbf{Methods:}
Scoping review following PRISMA-ScR \citep{Tricco2018} and JBI guidelines
\citep{Peters2020}. Searches were conducted in PubMed/MEDLINE, Scopus,
IEEE~Xplore, ACM~Digital~Library, Web~of~Science, and arXiv. The corpus
comprised [N]~articles selected from [N]~records, organised across four
thematic axes: (A)~clinical NLP and biomedical entity extraction;
(B)~automated clinical documentation and ASR; (C)~physician--patient
communication; (D)~vocal biomarkers and multimodal analysis.
%% TODO: replace [N] with final counts after search completion

\textbf{Results:}
Included studies document advances in clinical NLP (transformer-based models,
BERT leading named-entity recognition) and automated transcription (word error
rate $\approx$9\% with Whisper/PyAnnote). Patient-centred communication
demonstrates measurable impact on objective health outcomes. Communicational
markers --- self-repair, engagement, emotional support, and barrier indicators
--- are identified as adherence predictors. Multimodal acoustic--semantic
analysis remains nascent; most approaches focus on documentary automation
rather than treating clinical interaction as an analytic object.

\textbf{Conclusions:}
A consistent gap is confirmed: few studies operationalise clinical interactions
as measurable speech-derived variables, particularly communicational and
behavioural markers extracted from anamnesis. This review supports the design
of prospective studies integrating speech analytics into structured clinical
communication analysis for adherence monitoring and care quality assessment.

\textbf{Registration:} PROSPERO [CRD number to be assigned]
\end{abstract}

%% ----------------------------------------------------------------
%% Research Highlights — MANDATORY for IJMI (3 items, max 85 chars each)
%% ----------------------------------------------------------------
\begin{highlights}
\item BERT-based models lead clinical NER; WER $\approx$9\% achievable
      with Whisper/PyAnnote pipelines
\item Communicational markers (self-repair, engagement, barrier indicators)
      predict therapeutic adherence
\item Consistent gap: clinical interactions rarely operationalised as
      measurable speech-derived variables
\end{highlights}

%% ----------------------------------------------------------------
%% Keywords
%% ----------------------------------------------------------------
\begin{keywords}
Speech analytics \sep
Natural language processing \sep
Clinical communication \sep
Automatic speech recognition \sep
Communicational markers \sep
Digital scribe
\end{keywords}

\maketitle

%% ================================================================
\section{Introduction}
\label{sec:intro}

Clinical communication constitutes one of the richest yet most underutilised
sources of health data. In public health systems, anamneses and therapeutic
orientations predominantly occur as unstructured verbal exchanges rarely
converted into analysable data \citep{Wang2018,Koleck2019}.

The quality of the professional--patient interaction has been shown to
influence adherence and objective clinical outcomes \citep{Kelley2014}.
Communication barriers --- including health literacy gaps, misaligned
expectations, and miscommunication patterns --- are associated with impaired
therapeutic adherence and poorer prognosis \citep{McCabe2018,Alkhamees2023}.

Speech analytics --- the computational analysis of spoken language --- offers a
pathway to extract structured, measurable indicators from verbal clinical
interactions. Recent advances in ASR, particularly transformer-based models
such as Whisper \citep{Radford2023}, alongside instruction-tuned large
language models (LLMs), have lowered barriers to deploying end-to-end NLP
pipelines in healthcare \citep{Klusty2025}. Nonetheless, most existing
approaches target diagnostic support or documentation automation rather than
treating the clinical interaction itself as an object of structured analysis
\citep{vanBuchem2021}. The development of a speech analytics protocol for
monitoring orthopedic rehabilitation in a public health system
\citep{BaenaNeto2026} exemplifies the translational potential of these
technologies and underscores the need for a consolidated evidence map to
support prospective study design.

This scoping review aims to map the evidence on speech analytics and NLP
applied to clinical communication, characterising the methodological
landscape and identifying gaps for future prospective studies.

%% ================================================================
\section{Methods}
\label{sec:methods}

\subsection{Study design}

This scoping review follows the methodological framework proposed by
\citet{Arksey2005}, refined by \citet{Levac2010} and the Joanna Briggs
Institute (JBI) \citep{Peters2020}. Reporting adheres to the Preferred
Reporting Items for Systematic Reviews and Meta-Analyses extension for
Scoping Reviews (PRISMA-ScR) \citep{Tricco2018}.

\subsection{Research questions}

\begin{itemize}
  \item \textbf{Q1:} What speech analytics approaches have been applied to
        clinical communication in the last five years?
  \item \textbf{Q2:} What NLP and ASR methods are used for automated
        transcription and analysis of clinical consultations?
  \item \textbf{Q3:} What communicational and behavioural markers have been
        identified or proposed in professional--patient interaction?
  \item \textbf{Q4:} What are the technical, ethical, and methodological
        challenges for automated analysis of clinical communication?
\end{itemize}

\subsection{Eligibility criteria}

\subsubsection{Inclusion criteria}
\begin{itemize}
  \item \textbf{IC1:} Study applies speech analytics, NLP, or ASR to a
        clinical context
  \item \textbf{IC2:} Full text available in English, Portuguese, or Spanish
  \item \textbf{IC3:} Analyses professional--patient communication or
        information extraction from clinical consultations
  \item \textbf{IC4:} Describes methods, results, or markers applicable to
        clinical interaction
  \item \textbf{IC5:} Published between January 2014 and 2026
\end{itemize}

\subsubsection{Exclusion criteria}
\begin{itemize}
  \item \textbf{EC1:} Publications prior to 2014
  \item \textbf{EC2:} Extended abstracts, tutorials, or white papers without
        empirical data
  \item \textbf{EC3:} Exclusively focused on voice-based diagnosis of physical
        pathologies without analysis of the communicational interaction
  \item \textbf{EC4:} Studies on non-clinical populations without explicitly
        stated methodological transferability
\end{itemize}

\subsection{Information sources and search strategy}

Searches were conducted in six electronic databases: PubMed/MEDLINE, Scopus,
IEEE~Xplore, ACM~Digital~Library, Web~of~Science, and arXiv. The reference
list of \citet{BaenaNeto2026} was also hand-searched for additional records.
Searches were performed between [DATE\_START] and [DATE\_END].
%% TODO [Murilo]: insert search dates

The combined Boolean search string below was adapted to the controlled
vocabulary of each database. Database-specific strings are provided as
supplementary material.

\begin{verbatim}
("speech analytics" OR "automatic speech recognition" OR "ASR"
OR "natural language processing" OR "NLP" OR "large language model"
OR "speech-to-text" OR "voice recognition" OR "vocal biomarkers"
OR "speaker diarization")
AND
("clinical communication" OR "doctor-patient" OR "physician-patient"
OR "anamnesis" OR "medical consultation" OR "health interaction"
OR "patient-clinician communication" OR "therapeutic adherence")
AND
("clinical notes" OR "SOAP notes" OR "medical transcription"
OR "digital scribe" OR "electronic health record"
OR "communicational markers" OR "behavioural markers")
\end{verbatim}

\subsection{Study selection}

Records retrieved from the six databases were imported into Rayyan
\citep{Ouzzani2016} for deduplication and screening. Titles and abstracts
were independently screened by two reviewers (F.E.S. and M.M.), with
inter-rater agreement assessed using Cohen's kappa ($\kappa$). Disagreements
were resolved by a third reviewer (L.K.T.) or by the corresponding author.
Full texts of potentially eligible records were then retrieved and assessed
against the eligibility criteria. The selection process is documented in a
PRISMA-ScR flow diagram (Fig.~\ref{fig:prisma}).
%% TODO [Lucas]: insert final kappa value once calculated
%% TODO [Murilo]: insert final n per stage after all database searches complete

\subsection{Data charting}

Data were extracted using a standardised charting form including: author(s),
year, country, study design, clinical setting, population characteristics,
technology stack (ASR model, NLP framework, LLM), clinical variables
extracted, outcomes reported, ethical framework, and limitations.

\subsection{Synthesis}

A narrative synthesis was conducted. Findings are organised around four
thematic axes: (A)~clinical NLP and biomedical entity extraction,
(B)~automated documentation and ASR, (C)~physician--patient communication,
and (D)~vocal biomarkers and multimodal analysis.

%% ================================================================
\section{Results}
\label{sec:results}

% ---------------------------------------------------------------
\subsection{Study selection}
% ---------------------------------------------------------------

%% TODO [Murilo]: insert PRISMA-ScR flow figure and fill all [N] values
%%   after all database searches are complete and deduplicated in Rayyan.
%%   Template: Records identified (PubMed=[N], Scopus=[N], IEEE=[N],
%%   ACM=[N], WoS=[N], arXiv=[N]); duplicates removed=[N];
%%   screened=[N]; excluded on title/abstract=[N]; full-text assessed=[N];
%%   excluded on full-text (with reasons)=[N]; included=[N].

Database searches retrieved [N]~records in total
(PubMed/MEDLINE:~[N]; Scopus:~[N]; IEEE~Xplore:~[N];
ACM~Digital~Library:~[N]; Web~of~Science:~[N]; arXiv:~[N]).
After removal of [N]~duplicates, [N]~records were screened by title and
abstract. Inter-rater agreement was $\kappa = $ [VALUE] (Fabricio and
Murilo as primary reviewers). Following full-text assessment of [N]~records,
[N]~articles met all eligibility criteria and were included in the synthesis.
A PRISMA-ScR flow diagram is presented in Fig.~\ref{fig:prisma}.

\begin{figure}[h]
  \centering
  %% TODO [Murilo]: insert PRISMA-ScR flow diagram image here
  %% \includegraphics[width=\linewidth]{prisma_flow.pdf}
  \caption{PRISMA-ScR flow diagram of the study selection process.}
  \label{fig:prisma}
\end{figure}

% ---------------------------------------------------------------
\subsection{Characteristics of included studies}
% ---------------------------------------------------------------

%% TODO [Murilo]: populate the table below with all included studies.
%%   One row per study. Columns: Code | First author, year | Type |
%%   Axis | Technology | Main outcome/metric | Country

The [N]~included studies span [YEAR\_RANGE]. Table~\ref{tab:studies}
summarises their key characteristics.

\begin{longtable}{p{0.8cm} p{2.8cm} p{1.6cm} p{0.8cm} p{2.5cm} p{3.0cm}}
\caption{Characteristics of included studies.}
\label{tab:studies}\\
\toprule
\textbf{Code} & \textbf{Author, year} & \textbf{Type} &
\textbf{Axis} & \textbf{Technology} & \textbf{Main outcome} \\
\midrule
\endfirsthead
\multicolumn{6}{l}{\small\textit{(continued)}}\\
\toprule
\textbf{Code} & \textbf{Author, year} & \textbf{Type} &
\textbf{Axis} & \textbf{Technology} & \textbf{Main outcome} \\
\midrule
\endhead
\bottomrule
\endfoot
%% TODO [Murilo]: add one row per included study, e.g.:
%% S01 & Bazoge et al., 2023 & Syst.\ review & A & BERT/transformers
%%     & NER F1, phenotyping \\
%% S02 & Navarro et al., 2023 & Syst.\ review & A & BERT, BioBERT
%%     & NER entities (problems, treatments) \\
[N] & [Author, year] & [Type] & [X] & [Technology] & [Outcome] \\
\end{longtable}

% ---------------------------------------------------------------
\subsection{Axis A --- Clinical NLP and biomedical entity extraction}
\label{subsec:axisA}
% Responsible: Murilo Minutti
% ---------------------------------------------------------------

%% TODO [Murilo]: write the narrative synthesis for Axis A.
%%   Key studies to include: Bazoge2023, Navarro2023, Watabe2025,
%%   Sezgin2023, Koleck2019, Wang2018.
%%   Focus: transition from rule-based to transformer models,
%%   BERT/BioBERT leadership in NER, symptom extraction from
%%   free-text and patient-generated data, zero-shot approaches,
%%   performance metrics (F1 > 0.8 in Watabe2025), gaps in
%%   real-world deployment.
%%   Cite each claim with its source. Do not assert results not
%%   present in the extracted data.

[Axis A narrative — to be completed by Murilo]

% ---------------------------------------------------------------
\subsection{Axis B --- Automated clinical documentation and ASR}
\label{subsec:axisB}
% Responsible: Fabricio Evangelista dos Santos
% ---------------------------------------------------------------

%% TODO [Fabricio]: write the narrative synthesis for Axis B.
%%   Key studies to include: Krishna2021, Klusty2025, Abbasian2024,
%%   Aleksic2017, Tran2023, Schloss2020, vanBuchem2021.
%%   Focus: SOAP note generation, ASR pipelines (Whisper/PyAnnote,
%%   WER ~9%), speaker diarisation, digital scribe design,
%%   on-premise privacy requirements, LLM safety metrics
%%   (groundedness/RAG), challenges in colloquial speech.
%%   Cite each claim with its source. Do not assert results not
%%   present in the extracted data.

[Axis B narrative — to be completed by Fabricio]

% ---------------------------------------------------------------
\subsection{Axis C --- Physician--patient communication}
\label{subsec:axisC}
% Responsible: Lucas Kenso Terrazzin
% ---------------------------------------------------------------

%% TODO [Lucas]: write the narrative synthesis for Axis C.
%%   Key studies to include: Shao2025, Muscat2020, Kelley2014,
%%   Rucinski2022, McCabe2018, Alkhamees2023.
%%   Focus: DPCQ scale (Shao2025), impact of communication quality
%%   on objective health outcomes (Kelley2014 meta-analysis),
%%   self-repair as adherence predictor (McCabe2018), shared
%%   decision-making and health literacy (Muscat2020), intercultural
%%   barriers (Alkhamees2023), measurement gaps.
%%   Cite each claim with its source. Do not assert results not
%%   present in the extracted data.

[Axis C narrative — to be completed by Lucas]

% ---------------------------------------------------------------
\subsection{Axis D --- Vocal biomarkers and multimodal analysis}
\label{subsec:axisD}
% Responsible: Fabricio Evangelista dos Santos
% ---------------------------------------------------------------

%% TODO [Fabricio]: write the narrative synthesis for Axis D.
%%   Key studies to include: Fagherazzi2021, Albert2024, Imel2022.
%%   Focus: acoustic feature pipelines (prosody, voice quality),
%%   topic modelling from consultation audio (Imel2022),
%%   vocal biomarkers for health-state detection (Fagherazzi2021),
%%   gap in non-verbal/silence analysis (Albert2024),
%%   nascent state of multimodal acoustic+semantic integration.
%%   Cite each claim with its source. Do not assert results not
%%   present in the extracted data.

[Axis D narrative — to be completed by Fabricio]

%% ================================================================
\section{Discussion}
\label{sec:discussion}

%% TODO: complete after results synthesis

%% ================================================================
\section{Conclusion}
\label{sec:conclusion}

%% TODO: complete

%% ================================================================
\section*{Declarations}

\subsection*{Funding}
This study received no specific funding.
The article processing charge (APC) is covered under the Read and Publish
agreement between CAPES and Elsevier (2026--2028).

\subsection*{Conflict of interest}
The authors declare no conflict of interest.

\subsection*{Ethics approval}
Not applicable. This scoping review is based on publicly available
published literature and does not involve human participants.

\subsection*{Data availability}
The data charting form and full search strategies are available as
supplementary material upon acceptance.

\subsection*{AI use disclosure}
In accordance with Elsevier's Generative AI Policy: AI tools were used
solely for language support and literature searches. All intellectual
content, analysis, and conclusions are the sole responsibility of the
authors.

\printcredits

%% ================================================================
\bibliographystyle{cas-model2-names}
\bibliography{references}

\end{document}